\documentclass[conference]{IEEEtran}
\IEEEoverridecommandlockouts
\usepackage{cite}
\usepackage[T1]{fontenc}
\usepackage[utf8]{inputenc}
\usepackage{amsmath,amssymb,amsfonts}
\usepackage{algorithmic}
\usepackage{graphicx}
\usepackage{textcomp}
\usepackage{xcolor}
\usepackage{booktabs}
\usepackage{longtable}
\usepackage{array}
\usepackage{multirow}
\usepackage{url}
\def\BibTeX{{\rm B\kern-.05em{\sc i\kern-.025em b}\kern-.08em
    T\kern-.1667em\lower.7ex\hbox{E}\kern-.125emX}}

\begin{document}

\title{Engineering the Separation Principle: From Modular Architecture to Deployable Drone Navigation}

\author{\IEEEauthorblockN{Yusuf Saib}\\
\IEEEauthorblockA{\textit{Coybot Technology Corporation, Santa Clara, CA}\\
contact@coy.bot}}

\maketitle

\begin{abstract}
Our prior work established the separation principle for
vision-language drone navigation: VLMs should handle semantics while
dedicated modules handle geometry and safety. The resulting modular
pipeline achieved 1.04 m error on operational commands with 100\%
collision-free flight, but tied hover in aggregate (9.98 m vs.~9.50 m).
This paper reports the systematic engineering effort to close that
gap---and the deeper investigation that followed. Through eighteen
iterative refinements and 3,000+ closed-loop Isaac Sim trials, we first
improve the modular pipeline from 9.98 m to 8.15 m aggregate error,
beating hover for the first time with $\\geq$90\% collision-free flight. We
then construct Yonder, a 6.7-million-frame drone-perspective dataset
across 275 indoor environments, and use it to fine-tune open-vocabulary
detectors---achieving 9.7$\\times$ improvement in detection mAP (4.8\% $\\rightarrow$
46.7\%). However, four separate fine-tuning attempts across two detector
architectures fail to improve closed-loop navigation success over
zero-shot baselines. The root cause is a cross-simulator domain gap:
bounding box geometry calibrated to the training simulator (Habitat-Sim)
produces systematic localization errors in the evaluation simulator
(Isaac Sim), including negative-altitude goal predictions that cause the
drone to dive into the floor. With the domain gap identified and
patched, success rate converges at 23--25\% across all detector
variants---fine-tuned and zero-shot---revealing that detection is no
longer the binding constraint. Exploration is: 75\% of benchmark targets
are not visible from the drone's spawn position, and neither learned
exploration policies nor depth-based heuristics significantly improve
performance over random waypoints (+5.2 pp, p=0.168). We characterize
the remaining gap as requiring spatial reasoning and multi-step planning
capabilities that no component in the current modular architecture
addresses. All code, data, the Yonder dataset, and the complete
eighteen-iteration history are publicly available.
\end{abstract}

\\section{Introduction}\label{i.-introduction}

Our previous work [1] revealed a stark finding: across 10,200
closed-loop quadrotor trials evaluating 25 VLM architectures, no
end-to-end vision-language model reliably outperforms hovering in place.
The diagnosis was a \emph{metric gap}---VLMs achieve 0.83--0.91
directional accuracy but 6--10 m distance error on 4--12 m targets. The
prescription was the \emph{separation principle}: restrict VLMs to
semantic target identification and delegate spatial grounding to
dedicated depth and detection modules. The resulting modular pipeline
achieved 1.04 m mean error on operational commands with 100\%
collision-free flight.

However, the aggregate benchmark result told a less optimistic story.
Across the full 67-task suite, the hardened system scored 9.98
m---marginally worse than hover's 9.50 m, with the system winning on
only 21 of 51 individual tasks. This paper reports the engineering
effort to close that gap, and the deeper investigation into what
actually limits autonomous drone navigation.

The work proceeded in three phases, each driven by a hypothesis about
the binding constraint:

\textbf{Phase I (R3--R11): Pipeline engineering.} We hypothesized that
detection inconsistency, lack of temporal memory, and absence of
compound command support explained the aggregate gap. Six targeted
improvements---instruction decomposition, domain-specific detection
fine-tuning, vocabulary-constrained grounding, spatial semantic memory,
active yaw-based perception, and environment-specific safety
profiles---reduced aggregate error to 8.15 m, beating hover for the
first time with $\\geq$90\% collision-free flight.

\textbf{Phase II (R12--R17): Large-scale detection training.} We
hypothesized that more training data would further close the gap. We
constructed Yonder, a 6.7-million-frame dataset across 275 indoor
environments, and fine-tuned two detector architectures (OWL-ViT v2 and
Grounding DINO). Detection mAP improved 9.7$\\times$ (4.8\% $\\rightarrow$ 46.7\%). But
across four separate fine-tuning attempts, closed-loop navigation
success never exceeded the zero-shot baseline. The root cause was a
cross-simulator domain gap: bounding box geometry calibrated to
Habitat-Sim produced systematic localization errors in Isaac Sim.

\textbf{Phase III (R18): Exploration.} With detection ruled out as the
bottleneck, we hypothesized that active exploration would help---75\% of
targets are not visible from spawn. Both a learned exploration policy
and a depth-based heuristic produced marginal improvement (+5.2 pp,
p=0.168). The tiers that score 0\% (spatial reasoning, negation,
occluded targets, multi-step commands) require capabilities that no
component in the current architecture addresses.

Our contributions are:

\textbf{1) Pipeline engineering} that reduces aggregate error from 9.98
m to 8.15 m through six targeted improvements, validated across eleven
iterations and 2,000+ closed-loop trials.

\textbf{2) Yonder}, a 6.7-million-frame drone-perspective dataset
across 275 indoor environments with semantic segmentation, depth,
stereo, and LiDAR---the largest public dataset for drone indoor
navigation training.

\textbf{3) A negative result with broad implications}: 9.7$\\times$ improvement
in offline detection mAP produces zero improvement in closed-loop
navigation, due to a cross-simulator domain gap that systematically
corrupts 3D localization. This finding generalizes to any system that
trains on one simulator and evaluates on another.

\textbf{4) Bottleneck characterization} showing that after detection is
solved, the binding constraints are spatial reasoning and multi-step
planning---capabilities that require architectural changes, not better
perception.

\textbf{5) A complete eighteen-iteration engineering history}
documenting every hypothesis, test, success, and failure---providing a
roadmap for practitioners building deployable VLM-based navigation
systems.

\textbf{6) Swarm-scale validation} showing the pipeline achieves 70\%
success rate at n=4 drones versus 5\% for zero-shot VLM (p=0.002, 10
seeds, Bonferroni-corrected).

\\section{Related Work}\label{ii.-related-work}

Modular Navigation Pipelines.

The separation of semantic understanding from metric grounding has
precedent in classical robotics. SayCan [2] and SayNav [3] use
LLMs for high-level planning with skill primitives. LM-Nav [4]
constructs navigation graphs from CLIP and GPT-3 landmarks. VLMaps
[5] builds spatial language grounding from CLIP feature grids---an
approach our semantic memory map extends to the drone domain with 3D
backprojection and temporal accumulation. Our prior work [1]
established the separation principle specifically for drone navigation;
this paper demonstrates the engineering required to make it work in
aggregate, and documents where it reaches its limits.

Open-Vocabulary Detection for Navigation.

Grounding DINO [6] and OWL-ViT [7] enable text-conditioned
object detection without fixed category sets. Their zero-shot
performance is impressive on internet imagery but degrades on
drone-perspective views [1]. We demonstrate that domain-specific
fine-tuning dramatically improves offline detection metrics (9.7$\\times$ mAP
improvement) but does not transfer to closed-loop navigation when
training and evaluation use different simulators---a finding with
implications for the growing body of work on simulation-trained
perception for robotics.

Spatial Memory for Embodied Agents.

VLMaps [5] accumulates CLIP features in a bird's-eye-view grid.
SplaTAM [8] builds dense 3D Gaussian maps for SLAM. NaVid [9]
maintains detected object positions across video frames. Our semantic
memory map is deliberately simple---a sparse landmark store with EMA
updates and confidence-based filtering---because it must run alongside
detection and planning within a 200 ms per-cycle budget on edge
hardware.

Simulation-to-Simulation Transfer.

The sim-to-real gap is well-documented [15]. Our work identifies a
less-discussed problem: the \emph{sim-to-sim} gap. Models trained on
Habitat-Sim renders fail on Isaac Sim renders despite both being
physics-based indoor simulators. This manifests as broken confidence
calibration (OWL-ViT v2), negative-altitude goal predictions (GDINO
fine-tuned), and lateral localization errors---each requiring a
different diagnosis and fix.

\section{Baseline System And Limitations}\label{iii.-baseline-system-and-limitations}

Our starting point is the hardened Track A pipeline from [1]:

\textbf{Architecture.} An IRIS quadrotor in Isaac Sim with Mellinger
controller, 640$\\times$480 RGB-D camera (60° FOV), and four modules: VLM target
selector with Grounding DINO, Depth Anything V2 for metric depth, depth
backprojection with EMA smoothing, and classical planner with
depth-conditioned safety.

\textbf{Performance.} On operational commands: 1.04 m error, 100\%
collision-free. On the full benchmark: 9.98 m aggregate (vs.~9.50 m
hover), 100\% collision-free, 21/51 tasks beating hover.

\textbf{Limitations.} (L1) Detection inconsistency: hospital precision
0.28, office 0.20. (L2) No temporal memory: missed detection = wasted
cycle. (L3) No compound command support.

\section{Phase I: Pipeline Engineering (R3--R11)}\label{iv.-phase-i-pipeline-engineering-r3r11}

\subsection{Six Targeted Improvements}\label{a.-six-targeted-improvements}

We address each limitation through six improvements, each motivated by a
specific failure mode:

\emph{Instruction Decomposer (L3):} Qwen2.5-3B with QLoRA on 10K
synthetic pairs, with a vocabulary constraint layer mapping sub-goals to
GDINO-detectable terms (100\% coverage).

\emph{Domain-Specific Detection (L1):} Grounding DINO fine-tuned on
7,944 Isaac Sim frames. Hospital recall tripled (0.18 $\\rightarrow$ 0.54).

\emph{Failure Detector (safety):} ResNet-18 binary classifier. Three
training iterations required: v2 output NaN (synthetic data gap), v3
miscalibrated (97.9\% veto rate), v4 achieved 5.7\% veto rate with 98\%
collision-free flight.

\emph{Spatial Semantic Memory Map (L2):} Persistent 3D landmark store
with EMA updates. Merge radius 1.5 m, scan confirmation threshold 1.
Raised detection confirmation from 8.6\% to 81.5\%.

\emph{Active Yaw-Based Perception (L2):} 4-heading yaw scan before
navigation. An initial fly-to-offset scan improved detection by 0.14 m
but caused 18 pp collision-free drop; yaw-only scan eliminated collision
risk (99\% CF during scanning).

\emph{Environment-Specific Safety Profiles (safety):} Reactive collision
avoidance, hover threshold 1.5 m, max step 4.0 m for warehouse. Reduced
warehouse collisions by 62\%.

\subsection{Phase I Iteration History}\label{b.-phase-i-iteration-history}

\textbf{TABLE I:} Phase I iterations. All closed-loop Isaac Sim except
R3.

{\def\LTcaptype{none} % do not increment counter
\begin{longtable}[]{@{}
  >{\raggedright\arraybackslash}p{(\linewidth - 8\tabcolsep) * \real{0.2000}}
  >{\raggedright\arraybackslash}p{(\linewidth - 8\tabcolsep) * \real{0.2000}}
  >{\raggedright\arraybackslash}p{(\linewidth - 8\tabcolsep) * \real{0.2000}}
  >{\raggedright\arraybackslash}p{(\linewidth - 8\tabcolsep) * \real{0.2000}}
  >{\raggedright\arraybackslash}p{(\linewidth - 8\tabcolsep) * \real{0.2000}}@{}}
\toprule\noalign{}
\begin{minipage}[b]{\linewidth}\raggedright
Iteration
\end{minipage} & \begin{minipage}[b]{\linewidth}\raggedright
Error
\end{minipage} & \begin{minipage}[b]{\linewidth}\raggedright
CF\%
\end{minipage} & \begin{minipage}[b]{\linewidth}\raggedright
Key Change
\end{minipage} & \begin{minipage}[b]{\linewidth}\raggedright
Key Finding
\end{minipage} \\
\midrule\noalign{}
\endhead
\bottomrule\noalign{}
\endlastfoot
Paper 1 & 9.98 m & 100\% & Baseline & \\
R3 (mock) & 2.27 m & 100\% & +Decomposer, +FD v2 & Mock is 6.5$\\times$
optimistic \\
R4 (first CL) & 14.68 m & 88\% & First closed-loop & Decomposer hurts;
FD outputs NaN \\
R6 & 14.29 m & 88\% & +Vocab constraint, +FD v3 & FD still
miscalibrated \\
R7 & 8.49 m & 98\% & +GDINO fine-tune, +FD v4 & \textbf{First time
beating hover} \\
R8 & 8.93 m & 78\% & +Map, +scan, +routing & Scan causes collisions;
routing hurts \\
R9 & 8.97 m & 85\% & Yaw scan, relaxed map & Most tasks helped
(39/65) \\
R10 & 8.15 m & 94\% & +Warehouse safety & Best error + CF combination \\
R11 & 8.73 m & 90\% & Bug fix & Honest numbers \\
\end{longtable}
}

\subsection{Phase I Engineering Lessons}\label{c.-phase-i-engineering-lessons}

\emph{Mock evaluations are unreliable.} The 6.5$\\times$ gap between offline
(2.27 m) and closed-loop (14.68 m) is the most important methodological
finding. The decomposer's contribution reversed sign between modalities.

\emph{Synthetic data produces domain gaps.} The failure detector trained
on synthetic data output NaN on all Isaac Sim frames. The waypoint
prediction head trained on synthetic grids achieved 100\% collision
rate.

\emph{Complexity that hurts.} The spatial selector, frontier explorer,
and command router all degraded performance versus simpler alternatives.
Each was removed after ablation.

\\subsection{Phase I Final Results}\label{d.-phase-i-final-results}

\textbf{TABLE II:} Phase I aggregate results (R10, 198 trials).

{\def\LTcaptype{none} % do not increment counter
\begin{longtable}[]{@{}
  >{\raggedright\arraybackslash}p{(\linewidth - 12\tabcolsep) * \real{0.1429}}
  >{\raggedright\arraybackslash}p{(\linewidth - 12\tabcolsep) * \real{0.1429}}
  >{\raggedright\arraybackslash}p{(\linewidth - 12\tabcolsep) * \real{0.1429}}
  >{\raggedright\arraybackslash}p{(\linewidth - 12\tabcolsep) * \real{0.1429}}
  >{\raggedright\arraybackslash}p{(\linewidth - 12\tabcolsep) * \real{0.1429}}
  >{\raggedright\arraybackslash}p{(\linewidth - 12\tabcolsep) * \real{0.1429}}
  >{\raggedright\arraybackslash}p{(\linewidth - 12\tabcolsep) * \real{0.1429}}@{}}
\toprule\noalign{}
\begin{minipage}[b]{\linewidth}\raggedright
Architecture
\end{minipage} & \begin{minipage}[b]{\linewidth}\raggedright
Error
\end{minipage} & \begin{minipage}[b]{\linewidth}\raggedright
@1m
\end{minipage} & \begin{minipage}[b]{\linewidth}\raggedright
@3m
\end{minipage} & \begin{minipage}[b]{\linewidth}\raggedright
@5m
\end{minipage} & \begin{minipage}[b]{\linewidth}\raggedright
CF\%
\end{minipage} & \begin{minipage}[b]{\linewidth}\raggedright
Tasks \textgreater{} Hover
\end{minipage} \\
\midrule\noalign{}
\endhead
\bottomrule\noalign{}
\endlastfoot
Oracle & 0.15 m & 100\% & 100\% & 100\% & 100\% & 66/66 \\
Hover & 9.50 m & 6\% & 6\% & 16\% & 100\% & --- \\
Paper 1 & 9.98 m & 13\% & 19\% & 25\% & 100\% & 21/51 \\
\textbf{Phase I final} & \textbf{8.15 m} & \textbf{15\%} & \textbf{26\%}
& \textbf{34\%} & \textbf{94\%} & \textbf{32/66} \\
\end{longtable}
}

Warehouse: 5.98 m (89\% CF). Hospital: 8.31 m (100\% CF). Office: 12.07
m (100\% CF).

Error decomposition revealed: when detection succeeds, the pipeline
achieves 0.24 m (near-oracle). The dominant failure mode is detection
failure causing hover (30.8\% of trials)---targets behind walls, not
detection model inadequacy.

\section{Phase Ii: Large-Scale Detection Training (R12--R17)}\label{v.-phase-ii-large-scale-detection-training-r12r17}

The Phase I error decomposition identified detection failure as the
binding constraint (30.8\% hover rate). We hypothesized that training a
detector on a larger, more diverse dataset would close this gap.

\\subsection{Yonder Dataset}\label{a.-zeroclaw-dataset}

We constructed Yonder by flying a simulated Holybro x500v2 drone
through 275 indoor 3D scenes in Habitat-Sim, collecting multi-modal
sensor data at dense waypoints. The statistics below reflect this
internal working snapshot, prior to the license-driven scene reduction
described in the Yonder dataset paper: the public release retains only
the 167 HSSD scenes and excludes the ReplicaCAD, Replica, and HM3D
scenes shown here, whose upstream licenses do not permit open
redistribution of derivative renders.

\textbf{TABLE III:} Yonder dataset statistics (internal snapshot used
for these experiments; see the Yonder dataset paper for the public
release's post-license-reduction figures).

{\def\LTcaptype{none} % do not increment counter
\begin{longtable}[]{@{}
  >{\raggedright\arraybackslash}p{(\linewidth - 2\tabcolsep) * \real{0.5000}}
  >{\raggedright\arraybackslash}p{(\linewidth - 2\tabcolsep) * \real{0.5000}}@{}}
\toprule\noalign{}
\begin{minipage}[b]{\linewidth}\raggedright
Metric
\end{minipage} & \begin{minipage}[b]{\linewidth}\raggedright
Value
\end{minipage} \\
\midrule\noalign{}
\endhead
\bottomrule\noalign{}
\endlastfoot
Scenes & 275 (167 HSSD, 84 ReplicaCAD, 18 Replica, 6 HM3D) \\
Scenes with semantic annotations & 110 \\
Total waypoints & 556,490 \\
Total frames (waypoint $\\times$ 12 yaw) & 6,677,880 \\
Sensors per waypoint & Stereo RGB, depth, LiDAR 360°, landing camera,
up/down IR \\
Semantic categories & 405 \\
COCO annotations (from semantic re-rendering) & 32M bounding boxes \\
Total dataset size & 4.15 TB \\
Generation cost & \textasciitilde\$80 (Vast.ai RTX 4090 instances) \\
\end{longtable}
}

Semantic segmentation was not stored during initial generation (used
only for adaptive waypoint placement). We re-rendered semantic channels
for all 110 scenes with semantic meshes, generating COCO-format bounding
box annotations by projecting instance masks to 2D.

\\subsection{Detection Benchmarking}\label{b.-detection-benchmarking}

Zero-shot detection on Yonder test frames:

{\def\LTcaptype{none} % do not increment counter
\begin{longtable}[]{@{}llll@{}}
\toprule\noalign{}
Model & mAP@0.5 & mAP@0.75 & Inference \\
\midrule\noalign{}
\endhead
\bottomrule\noalign{}
\endlastfoot
OWL-ViT v2 (zero-shot) & 4.8\% & 2.7\% & 168 ms \\
Grounding DINO (zero-shot) & \textasciitilde5\% & --- & 75 ms \\
\end{longtable}
}

Both models perform poorly on drone-perspective imagery---consistent
with [1]'s finding that VLM-based perception is uncalibrated for
this domain.

\subsection{Four Fine-Tuning Attempts}\label{c.-four-fine-tuning-attempts}

\textbf{TABLE IV:} Detection fine-tuning attempts and closed-loop
outcomes.

{\def\LTcaptype{none} % do not increment counter
\begin{longtable}[]{@{}
  >{\raggedright\arraybackslash}p{(\linewidth - 10\tabcolsep) * \real{0.1667}}
  >{\raggedright\arraybackslash}p{(\linewidth - 10\tabcolsep) * \real{0.1667}}
  >{\raggedright\arraybackslash}p{(\linewidth - 10\tabcolsep) * \real{0.1667}}
  >{\raggedright\arraybackslash}p{(\linewidth - 10\tabcolsep) * \real{0.1667}}
  >{\raggedright\arraybackslash}p{(\linewidth - 10\tabcolsep) * \real{0.1667}}
  >{\raggedright\arraybackslash}p{(\linewidth - 10\tabcolsep) * \real{0.1667}}@{}}
\toprule\noalign{}
\begin{minipage}[b]{\linewidth}\raggedright
Attempt
\end{minipage} & \begin{minipage}[b]{\linewidth}\raggedright
Model
\end{minipage} & \begin{minipage}[b]{\linewidth}\raggedright
Training Data
\end{minipage} & \begin{minipage}[b]{\linewidth}\raggedright
mAP@0.5
\end{minipage} & \begin{minipage}[b]{\linewidth}\raggedright
Closed-Loop SR@5m
\end{minipage} & \begin{minipage}[b]{\linewidth}\raggedright
Failure Mode
\end{minipage} \\
\midrule\noalign{}
\endhead
\bottomrule\noalign{}
\endlastfoot
R12 OWL-ViT & OWL-ViT v2 & 29K frames, 90 scenes & --- & 11.6\% (5
seeds) & Confidence calibration collapse (blank images score 0.97) \\
R14 GDINO R6 & GDINO-tiny & 40K frames, 20 scenes & \textbf{46.7\%} &
23.5\% & Negative-Z goals from bbox height shift \\
R16 GDINO R6 & Same & Same & Same & 25.5\% (post Z-clamp) & Tied with
zero-shot after clamp \\
R17 GDINO R7 & GDINO-tiny & 40K Yonder + 1.5K Isaac Sim & 45.4\% &
25.5\% & Circular pseudo-GT; warehouse CF collapse \\
\end{longtable}
}

Zero-shot baselines:

{\def\LTcaptype{none} % do not increment counter
\begin{longtable}[]{@{}
  >{\raggedright\arraybackslash}p{(\linewidth - 4\tabcolsep) * \real{0.3333}}
  >{\raggedright\arraybackslash}p{(\linewidth - 4\tabcolsep) * \real{0.3333}}
  >{\raggedright\arraybackslash}p{(\linewidth - 4\tabcolsep) * \real{0.3333}}@{}}
\toprule\noalign{}
\begin{minipage}[b]{\linewidth}\raggedright
Config
\end{minipage} & \begin{minipage}[b]{\linewidth}\raggedright
SR@5m
\end{minipage} & \begin{minipage}[b]{\linewidth}\raggedright
Source
\end{minipage} \\
\midrule\noalign{}
\endhead
\bottomrule\noalign{}
\endlastfoot
ZS GDINO + classical (RESULTS13) & \textbf{32.9\%} {[}28.8, 37.2{]} &
Real Isaac Sim, 5 seeds, 465 trials \\
ZS GDINO + safety + Z-clamp (RESULTS16) & 23.5\% & Kinematic, 3 seeds,
153 trials \\
\end{longtable}
}

No fine-tuned model exceeds the zero-shot baseline in closed-loop
evaluation. The 9.7$\\times$ mAP improvement does not transfer.

\subsection{Root Cause: Cross-Simulator Domain Gap}\label{d.-root-cause-cross-simulator-domain-gap}

Fine-tuning on Habitat-Sim data calibrates bounding box centers to
Habitat-Sim geometry. In Isaac Sim's different camera model and depth
scale, these predictions produce systematic 3D localization errors:

\emph{Negative-Z goals.} Floor-level objects (forklifts, pallet jacks)
have bounding box centers lower in frame in Habitat-Sim than in Isaac
Sim. Depth backprojection converts these to negative-Z world
coordinates, causing the drone to dive toward the floor. Fine-tuned
GDINO produces negative-Z goals 6$\\times$ more often than zero-shot (24\% vs
4\% of warehouse trials).

\emph{Confidence calibration collapse.} OWL-ViT v2 fine-tuning shifted
the sigmoid head to extreme values---blank images scored 0.97. The
detection threshold became meaningless (all predictions above any
threshold).

\emph{Lateral localization shift.} Even with Z-floor clamping,
fine-tuned bounding box centers send the drone into shelving units
rather than through aisle gaps.

A Z-floor clamp (\texttt{goal\_z\ =\ max(goal\_z,\ 0.5)}) eliminates the
floor-diving but does not recover SR@5m: the improvement is absorbed by
other localization errors. Adding 10\% Isaac Sim pseudo-labeled frames
to training (R17) failed because pseudo-labels from zero-shot GDINO
create a circular training signal.

\\subsection{The Detection Plateau}\label{e.-the-detection-plateau}

With the domain gap identified, SR@5m converges at 23--25\% across all
detector configurations (fine-tuned R6, fine-tuned R7, zero-shot, with
and without safety profiles). This convergence demonstrates that
\textbf{detection quality is no longer the binding constraint.} The
pipeline succeeds on targets it can see and fails on targets it
cannot---regardless of how well the detector performs on visible
targets.

\section{Phase Iii: Exploration (R18)}\label{vi.-phase-iii-exploration-r18}

With detection ruled out, we hypothesized that active exploration would
address the 75\% of targets not visible from spawn.

\subsection{Learned Exploration Policy}\label{a.-learned-exploration-policy}

We attempted to train a ResNet18-based heading predictor on Isaac Sim
depth panoramas with CLIP text embeddings. The training failed:
validation accuracy (6.9\%) was below random chance (8.3\%) on 4,704
samples across 3 environments---insufficient data to learn a
generalizable 12-class heading classifier.

The failure of the Phase 3 learned policy (RESULTS13: -21 pp SR) had
three diagnosed causes: (1) random text embeddings instead of real CLIP
during training, (2) 12-frame panorama capture overhead doubling
wall-clock time, and (3) Habitat-Sim depth used for training, Isaac Sim
depth for evaluation.

\subsection{Depth Heuristic Explorer}\label{b.-depth-heuristic-explorer}

As a fallback, we implemented a zero-training depth heuristic: at each
exploration step, fly 2 m toward the yaw heading with highest combined
mean and 90th-percentile depth (most open space). Maximum 3 exploration
steps per trial. Reuses cached yaw-scan depth frames (zero additional
capture overhead).

\\subsection{Exploration Results}\label{c.-exploration-results}

\textbf{TABLE V:} Exploration ablation (153 trials per condition, 3
seeds).

{\def\LTcaptype{none} % do not increment counter
\begin{longtable}[]{@{}llll@{}}
\toprule\noalign{}
Condition & SR@5m & CF\% & p vs no-exploration \\
\midrule\noalign{}
\endhead
\bottomrule\noalign{}
\endlastfoot
ZS GDINO + depth heuristic exploration & \textbf{24.8\%} & 86.9\% &
0.168 \\
ZS GDINO + no exploration & 19.6\% & 45.8\% & --- \\
\end{longtable}
}

The +5.2 pp improvement is not statistically significant (p=0.168,
Fisher exact test). Per-tier analysis reveals the fundamental
limitation:

\textbf{TABLE VI:} Per-tier success rate with depth heuristic
exploration.

{\def\LTcaptype{none} % do not increment counter
\begin{longtable}[]{@{}lll@{}}
\toprule\noalign{}
Tier & SR@5m & Status \\
\midrule\noalign{}
\endhead
\bottomrule\noalign{}
\endlastfoot
1 Stationary & 100\% & Solved \\
2 Cardinal & 100\% & Solved \\
11 Appearance & 66.7\% & Working \\
14 Exploration & 33.3\% & Partial \\
3 Named Near & 33.3\% & Partial \\
6 Visual & 33.3\% & Partial \\
13 Counting & 33.3\% & Partial \\
12 Semantic & 11.1\% & Marginal \\
9 Obstacle & 5.6\% & Failing \\
10 Occluded & 0\% & Unsolved \\
4 Named Far & 0\% & Unsolved \\
5 Spatial & 0\% & Unsolved \\
7 Reasoning & 0\% & Unsolved \\
8 Multi-leg & 0\% & Unsolved \\
\end{longtable}
}

The tiers scoring 0\% share a common requirement: they need capabilities
the pipeline does not have. Occluded targets require the drone to
navigate to vantage points it has never visited. Spatial reasoning
requires understanding ``closest to the entrance''---a spatial relation
between an object and a scene landmark. Negation requires selecting
``NOT the nearest'' from multiple detected instances. Multi-leg commands
require maintaining state across sequential sub-goals with adaptive time
budgeting. None of these are detection problems or exploration problems.
They are reasoning and planning problems.

\\section{Ablation Studies}\label{vii.-ablation-studies}

\subsection{Component Contributions (Phase I)}\label{a.-component-contributions-phase-i}

\textbf{TABLE VII:} Ablation results (closed-loop Isaac Sim, 54 trials
per configuration).

{\def\LTcaptype{none} % do not increment counter
\begin{longtable}[]{@{}llll@{}}
\toprule\noalign{}
Configuration & Error & CF\% & Key Effect \\
\midrule\noalign{}
\endhead
\bottomrule\noalign{}
\endlastfoot
Full pipeline & 8.97 m & 85\% & Baseline \\
No semantic map, no scan & 9.02 m & 93\% & Map + scan worth
\textasciitilde0.05 m \\
No scan (map only) & 10.10 m & 96\% & Scan worth \textasciitilde1.1 m \\
No routing (map + scan) & 9.19 m & 81\% & Routing \emph{hurts} by 0.55
m \\
\end{longtable}
}

\\subsection{Mock vs.~Closed-Loop}\label{b.-mock-vs.-closed-loop}

\textbf{TABLE VIII:} The evaluation gap.

{\def\LTcaptype{none} % do not increment counter
\begin{longtable}[]{@{}llll@{}}
\toprule\noalign{}
Metric & Mock & Closed-Loop & Gap \\
\midrule\noalign{}
\endhead
\bottomrule\noalign{}
\endlastfoot
Mean error & 2.27 m & 14.68 m & 6.5$\\times$ \\
@5m success & 100\% & 28\% & 3.6$\\times$ \\
Decomposer impact & +1.79 m (helps) & -6.09 m (hurts) &
\textbf{Reversed} \\
\end{longtable}
}

\\subsection{Swarm-Scale Validation}\label{c.-swarm-scale-validation}

We validated the pipeline at swarm scale using the cooperative framework
from [16], with 234 closed-loop trials across 10 seeds.

\textbf{TABLE IX:} Swarm results (n=4 drones, warehouse, 10 seeds).

{\def\LTcaptype{none} % do not increment counter
\begin{longtable}[]{@{}
  >{\raggedright\arraybackslash}p{(\linewidth - 6\tabcolsep) * \real{0.2500}}
  >{\raggedright\arraybackslash}p{(\linewidth - 6\tabcolsep) * \real{0.2500}}
  >{\raggedright\arraybackslash}p{(\linewidth - 6\tabcolsep) * \real{0.2500}}
  >{\raggedright\arraybackslash}p{(\linewidth - 6\tabcolsep) * \real{0.2500}}@{}}
\toprule\noalign{}
\begin{minipage}[b]{\linewidth}\raggedright
Architecture
\end{minipage} & \begin{minipage}[b]{\linewidth}\raggedright
Success Rate
\end{minipage} & \begin{minipage}[b]{\linewidth}\raggedright
95\% CI
\end{minipage} & \begin{minipage}[b]{\linewidth}\raggedright
p vs VLM
\end{minipage} \\
\midrule\noalign{}
\endhead
\bottomrule\noalign{}
\endlastfoot
Modular cooperative (FT GDINO + comms) & \textbf{72.5\%} & 65--80\% &
0.002 \\
Modular pipeline (FT GDINO, independent) & \textbf{70.0\%} & 62--75\% &
0.002 \\
VLM Qwen (zero-shot) & 5.0\% & 0--12\% & --- \\
Hover & 10.0\% & 0--20\% & --- \\
\end{longtable}
}

The fine-tuned pipeline outperforms zero-shot VLM by 65 pp at n=4
(p=0.002, Bonferroni-corrected). Cooperative detection sharing adds only
+2.5 pp (n.s.), confirming that per-drone perception quality dominates
over swarm coordination.

\subsection{Detection Fine-Tuning (Phase II)}\label{d.-detection-fine-tuning-phase-ii}

\textbf{TABLE X:} 2$\\times$2 ablation: fine-tuned vs zero-shot $\\times$ safety ON vs
OFF (153 trials per condition).

{\def\LTcaptype{none} % do not increment counter
\begin{longtable}[]{@{}lll@{}}
\toprule\noalign{}
& Safety ON & Safety OFF \\
\midrule\noalign{}
\endhead
\bottomrule\noalign{}
\endlastfoot
Fine-tuned GDINO & 23.5\% SR, 86.9\% CF & 20.4\% SR, 71.8\% CF \\
Zero-shot GDINO & 23.5\% SR, 86.9\% CF & 24.7\% SR, 80.0\% CF \\
\end{longtable}
}

No significant difference between any pair (all p \textgreater{} 0.25,
Bonferroni-corrected). Fine-tuning does not help in closed-loop.

\\section{Discussion}\label{viii.-discussion}

Three bottlenecks, sequentially resolved.

The iteration history reveals a cascading bottleneck structure. Each
time we resolved the dominant constraint, a new one became visible:

\emph{Bottleneck 1: Detection inconsistency (Phase I).} Zero-shot GDINO
failed on 30\% of targets. Pipeline engineering (fine-tuning, semantic
memory, active scanning) reduced aggregate error from 9.98 m to 8.15 m.
Detection rate in closed-loop reached 86--94\%.

\emph{Bottleneck 2: Cross-simulator domain gap (Phase II).} Fine-tuning
on Habitat-Sim produced 9.7$\\times$ better mAP but broke Isaac Sim
localization. Four attempts with two architectures produced the same
result: zero-shot GDINO remains optimal for closed-loop because it
doesn't carry domain-specific geometric biases. The 46.7\% mAP is real
but non-transferable.

\emph{Bottleneck 3: Reasoning and planning (Phase III).} With detection
solved (for visible targets) and exploration providing marginal benefit,
the remaining 0\% tiers require spatial reasoning, negation, and
multi-step planning. These are architectural gaps, not parameter tuning
opportunities.

The cross-simulator domain gap as a general finding.

The sim-to-sim transfer failure is our most surprising finding. Both
Habitat-Sim and Isaac Sim render physics-based indoor environments, yet
bounding box geometry trained on one does not transfer to the other. The
failure modes are specific and diagnosable---negative-Z predictions,
confidence collapse, lateral shifts---but they are invisible to offline
evaluation metrics (mAP improved 9.7$\\times$). This suggests that the commonly
assumed fungibility of simulation platforms for robotic perception
training may be more fragile than the field recognizes.

Implications for the ``just scale the data'' hypothesis.

Yonder's 6.7M frames across 275 environments represents substantial
training scale. The complete failure to improve closed-loop navigation
despite a 9.7$\\times$ offline improvement challenges the assumption that more
simulation data straightforwardly produces better robot performance. The
binding constraint shifted from data quantity to domain alignment---and
domain alignment cannot be resolved by adding more data from the wrong
domain.

What would actually help.

The unsolved tiers (spatial, reasoning, occluded, multi-leg) require
capabilities that the current modular pipeline cannot provide through
component-level improvements:

A \emph{spatial reasoning module} that understands relations between
objects and scene landmarks (``the shelf closest to the entrance'')
requires maintaining a model of the environment's semantic layout, not
just individual object positions.

A \emph{multi-step planner} that adaptively allocates time across
compound sub-goals and re-plans when intermediate goals fail requires
online replanning capabilities beyond the current look-fly-look
protocol.

An \emph{active exploration policy} that navigates toward likely target
locations based on scene structure (doorways lead to rooms containing
beds) requires either learned scene priors from many environments or a
topological map that the current system does not build.

These are architectural changes, not training data improvements. They
represent the next phase of work beyond the separation principle's
current formulation.

Implications for defense applications.

The pipeline reliably handles concrete, visible targets in structured
environments: warehouse forklift navigation at 95\% success (R12, single
seed), operational commands at 1.47 m error, and 100\% collision-free
flight in hospital and office. The system's failure mode (hover when
uncertain) is operationally correct. The modular architecture supports
NDAA-compliant component selection and component-level V\&V. Extension
to n=4 swarms achieves 70\% success rate with the pipeline operating
independently per drone---meaning the system functions under
communications denial.

\\section{Limitations}\label{ix.-limitations}

\begin{enumerate}
\def\labelenumi{(\arabic{enumi})}
\tightlist
\item
  All results are simulation-only; real-world transfer is unvalidated.
  (2) Indoor environments only. (3) The cross-simulator domain gap
  (Habitat-Sim $\\rightarrow$ Isaac Sim) may be specific to these two platforms and
  not generalize. (4) Yonder annotations use category-level labels,
  not instance-level; instance discrimination relies on spatial
  position, not visual identity. (5) The exploration policy evaluation
  used only 3 seeds (153 trials); the p=0.168 non-significance may be a
  power issue rather than a true null effect. (6) Phase I results (8.15
  m, 94\% CF) were obtained with fine-tuned GDINO; the zero-shot GDINO +
  safety profile + Z-clamp configuration achieves 23.5\% SR@5m on a
  harder 51-task benchmark not directly comparable to the Phase I
  66-task benchmark. (7) Swarm collision-free rates (12\%) are low
  because safety profiles do not account for inter-drone collisions. (8)
  The stall-based collision detector counts conservative stops as
  collisions, potentially underreporting true safety.
\end{enumerate}

\\section{Conclusion}\label{x.-conclusion}

Through eighteen iterative refinements and 3,000+ closed-loop Isaac Sim
trials, we systematically investigated three hypothesized bottlenecks
for autonomous drone navigation:

\emph{Detection} was the first bottleneck. Pipeline engineering (Phase
I) reduced aggregate error from 9.98 m to 8.15 m, beating hover for the
first time. A 6.7-million-frame dataset (Yonder) and four fine-tuning
attempts achieved 9.7$\\times$ offline mAP improvement but zero closed-loop
improvement due to cross-simulator domain gap.

\emph{Exploration} was the second hypothesis. Neither a learned policy
nor a depth heuristic significantly improved success rate (+5.2 pp,
p=0.168) over random waypoints.

\emph{Reasoning and planning} is the actual remaining bottleneck. The
tiers scoring 0\%---spatial reasoning, negation, multi-step commands,
occluded targets---require capabilities that the modular
perception-planning pipeline cannot provide through better detection,
more training data, or exploration heuristics.

The iteration history is itself a contribution. Mock evaluation
overestimated performance by 6.5$\\times$. Component-level mAP improvement did
not predict system-level improvement. Four fine-tuning attempts across
two architectures produced the same negative result. A sophisticated
routing system performed worse than the simple nearest-in-map default.
Each finding required closed-loop testing to discover.

The separation principle remains sound: VLMs for semantics, dedicated
modules for geometry and safety. But the principle has a ceiling.
Crossing it requires not better perception but better
reasoning---understanding spatial relations, planning multi-step
strategies, and navigating to places the drone has never been. We
release all code, all eighteen iterations of experimental data, the
Yonder dataset, and the complete pipeline to enable others to start
from where we stopped.

\begin{thebibliography}{99}
\bibitem{ref1} Y. Saib, ``Closing the Metric Gap: From Diagnosis to Solution in Vision-Language Drone Navigation,'' 2026.
\bibitem{ref2} M. Ahn et al., ``SayCan: Grounding Language in Robotic Affordances,'' arXiv:2204.01691, 2022.
\bibitem{ref3} K. Rajvanshi et al., ``SayNav: Grounding LLMs for Dynamic Planning,'' arXiv:2309.04077, 2023.
\bibitem{ref4} D. Shah et al., ``LM-Nav: Robotic Navigation with Large Pre-Trained Models,'' CoRL, 2023.
\bibitem{ref5} H. Huang et al., ``VLMaps: Visual Language Maps for Robot Navigation,'' IJRR, 2025.
\bibitem{ref6} S. Liu et al., ``Grounding DINO: Marrying DINO with Grounded Pre-Training for Open-Set Object Detection,'' ECCV, 2024.
\bibitem{ref7} M. Minderer et al., ``Scaling Open-Vocabulary Object Detection,'' NeurIPS, 2023.
\bibitem{ref8} N. Keetha et al., ``SplaTAM: Splat, Track \& Map 3D Gaussians for Dense RGB-D SLAM,'' CVPR, 2024.
\bibitem{ref9} J. Zhang et al., ``NaVid: Video-based VLM for VLN,'' RSS, 2024.
\bibitem{ref10} C. Stachniss et al., ``Information-Gain-Based Exploration Using Rao-Blackwellized Particle Filters,'' RSS, 2005.
\bibitem{ref11} A. Bircher et al., ``Receding Horizon Next-Best-View Planner for 3D Exploration,'' ICRA, 2016.
\bibitem{ref12} D. Simon et al., ``MonoNav: MAV Navigation via Monocular Depth Estimation,'' ISER, 2023.
\bibitem{ref13} NVIDIA, ``Isaac Sim: Robotics Simulation and Synthetic Data,'' 2024.
\bibitem{ref14} M. Jacinto et al., ``Pegasus Simulator: An Isaac Sim Framework for Multiple Aerial Vehicles Simulation,'' ICUAS, 2024.
\bibitem{ref15} J. Tobin et al., ``Domain Randomization for Transferring Deep Neural Networks from Simulation to the Real World,'' IROS, 2017.
\bibitem{ref16} Y. Saib, ``Scaling the Separation Principle: Sensing Requirements for 1000-Drone Swarms in Urban and Natural Environments,'' 2026.
\end{thebibliography}

\end{document}
